% =============================================================================
% APPENDIX EEE: METHOD-DESIGN: Model Design Workflow
% =============================================================================
% Module: BCM2_METHOD_DESIGN
% Category: METHOD
% Language: English
% Version: 1.0 (January 2026)
% Status: Draft
%
% This appendix defines the systematic workflow for designing behavioral models
% from the EBF framework. It provides entry points for both theory-driven and
% practice-driven approaches, with the 3+1 Choice Architecture for decisions.
%
% =============================================================================

% =============================================================================
% CROSS-REFERENCE STRUCTURE
% =============================================================================
%
% CHAPTER LINKAGE (Required):
% - Primary Chapter: Chapter 4: From Framework to Model
% - Secondary Chapters: Chapter 10 (Applications), Chapter 18 (Meta-Theory)
%
% DEPENDENCIES (what this appendix REQUIRES):
% - Appendix REG (METHOD-REGISTRY): Model archive for option selection
% - Appendix EST (CORE-WHERE): Parameter sources
% - Appendix CTW (CORE-WHEN): Context specification
% - Appendix WAT (CORE-WHAT): Utility dimensions
% - Appendix HOW (CORE-HOW): Complementarity specification
%
% DEPENDENTS (what REQUIRES this appendix):
% - All DOMAIN- appendices: Use this workflow for domain-specific models
% - All PREDICT- appendices: Models designed via this workflow
% - Appendix REG (METHOD-REGISTRY): Stores models created via this workflow
%
% =============================================================================

\begin{tcolorbox}[colback=purple!5!white, colframe=purple!75!black,
    title=Chapter Linkage]

\textbf{Primary Chapter:} Chapter 4: From Framework to Model
\begin{itemize}[nosep]
\item Section 4.1: Framework vs. Model Distinction $\leftrightarrow$ Section \ref{sec:eee-fundamental}
\item Section 4.2: Model Design Process $\leftrightarrow$ Section \ref{sec:eee-workflow}
\item Section 4.3: Entry Points $\leftrightarrow$ Section \ref{sec:eee-entry}
\end{itemize}

\textbf{Secondary Chapters:}
\begin{itemize}[nosep]
\item Chapter 10: Applications --- uses workflow for domain models
\item Chapter 18: Meta-Theory --- theoretical foundations of design
\end{itemize}

\textbf{Reading Guide:}
\begin{itemize}[nosep]
\item For conceptual overview: Read Chapter 4 first
\item For formal workflow: This appendix provides complete specification
\item For model storage: See Appendix REG (METHOD-REGISTRY)
\end{itemize}

\end{tcolorbox}

\chapter{METHOD-DESIGN: Model Design Workflow}
\label{app:method-design}

\begin{abstract}
This appendix defines the systematic workflow for designing testable behavioral
models from the Evidence-Based Framework (EBF). The workflow supports two entry
points---theory-driven and practice-driven (Gestaltungsanlass)---which converge
through iterative Scope-Context specification. Central to the workflow is the
3+1 Choice Architecture (Axiom MD-1), which provides three curated options plus
one self-defined option at each decision point. Options are drawn from the Model
Registry (Appendix REG) via similarity search (Axiom MD-2).
\end{abstract}

\begin{tcolorbox}[colback=blue!5!white, colframe=blue!75!black,
    title=Quick Reference: Model Design Workflow]
\textbf{Appendix:} DSN (METHOD)


\textbf{Core Question:} How do we systematically design a testable model from EBF?

\textbf{Key Principle:}
\begin{equation}
\text{Framework (EBF)} \xrightarrow{\text{Workflow}} \text{Model } M(\Psi, C, \gamma, \Theta)
\end{equation}

\textbf{Key Concepts:}
\begin{itemize}[nosep]
\item \textbf{Recursive 3+1}: Workflow selection itself follows 3+1 (Geführt/Schnell/Template/Custom)
\item \textbf{3+1 Choice Architecture}: 3 curated options + 1 self-defined at each decision
\item \textbf{Entry Points}: Theory-driven vs. Gestaltungsanlass (practice-driven)
\item \textbf{Convergence}: Iterative Scope $\leftrightarrow$ Context stabilization
\item \textbf{Registry Integration}: Options from validated historical models (FFF)
\end{itemize}

\textbf{Cross-References:} FFF (Registry), GGG (Configurator), BBB (Parameters), V (Context), C (Utility), B (Complementarity), CAL (LLMMC Estimation)
\end{tcolorbox}

% -----------------------------------------------------------------------------
% APPENDIX SCOPE BOX
% -----------------------------------------------------------------------------

\begin{tcolorbox}[
    colback=orange!5!white,
    colframe=orange!75!black,
    title={\textbf{Appendix Scope: Ziel / In-Scope / Out-of-Scope / Constraints}},
    fonttitle=\bfseries
]

\textbf{Ziel (Objective):}
\begin{itemize}[nosep]
    \item Aspect
\end{itemize}

\textbf{In-Scope:}
\begin{itemize}[nosep]
    \item The Fundamental Question
    \item Entry Points
    \item Axioms and Formal Definitions
    \item The Complete Workflow
    \item Integration with Other Appendices
\end{itemize}

\textbf{Out-of-Scope (delegiert an andere Appendices):}
\begin{itemize}[nosep]
    \item METHOD-REGISTRY $\rightarrow$ Appendix REG
    \item CORE-WHERE $\rightarrow$ Appendix EST
    \item CORE-WHEN $\rightarrow$ Appendix CTW
    \item CORE-WHAT $\rightarrow$ Appendix WAT
    \item CORE-HOW $\rightarrow$ Appendix HOW
\end{itemize}

\textbf{Constraints (Anwendungsgrenzen):}
\begin{itemize}[nosep]
    \item [TODO: Define when this appendix does NOT apply]
    \item [TODO: Define required prerequisites]
    \item [TODO: Define known limitations]
\end{itemize}

\textbf{Lieferobjekte (Deliverables):}
\begin{itemize}[nosep]
    \item 4 Table(s)
    \item 9 Equation(s)
    \item 9 Axiom(s)
    \item Worked Example(s)
\end{itemize}

\end{tcolorbox}




% =============================================================================
%                    PART 2: CORE CONTENT
% =============================================================================

%%%%%%%%%%%%%%%%%%%%%%%%%%%%%%%%%%%%%%%%%%%%%%%%%%%%%%%%%%%%%%%%%%%%%%%%%%%%%%%
\section{The Fundamental Question}
\label{sec:eee-fundamental}
%%%%%%%%%%%%%%%%%%%%%%%%%%%%%%%%%%%%%%%%%%%%%%%%%%%%%%%%%%%%%%%%%%%%%%%%%%%%%%%

\subsection{Why This Matters}

The Evidence-Based Framework (EBF) is a \textit{framework}---an organizational
structure that integrates 50+ years of behavioral economic evidence. However,
frameworks do not make predictions; \textit{models} do.

\begin{table}[htbp]
\centering
\caption{Framework vs. Model Distinction}
\label{tab:eee-framework-model}
\renewcommand{\arraystretch}{1.3}
\begin{tabular}{@{}lll@{}}
\toprule
\textbf{Aspect} & \textbf{Framework (EBF)} & \textbf{Model} \\
\midrule
Purpose & Organize evidence & Make predictions \\
Nature & Descriptive & Predictive \\
Statement & ``Behavior can be understood as...'' & ``Behavior will change by...'' \\
Testability & Not directly testable & Falsifiable \\
Scope & Universal & Context-specific \\
\bottomrule
\end{tabular}
\end{table}

The transition from framework to model requires a systematic workflow that
ensures rigor, consistency, and traceability.

\subsection{The Core Problem}

Without a defined workflow:
\begin{itemize}
\item Models are designed ad-hoc, without leveraging organizational learning
\item The same ``wheel'' is reinvented across projects
\item Model quality depends entirely on individual consultant expertise
\item No audit trail exists for model-based recommendations
\end{itemize}

\begin{tcolorbox}[colback=red!5!white, colframe=red!75!black,
    title=The Fundamental Question]
\textbf{How do we systematically transform EBF from an organizing framework into
a specific, testable, context-appropriate behavioral model?}
\end{tcolorbox}


%%%%%%%%%%%%%%%%%%%%%%%%%%%%%%%%%%%%%%%%%%%%%%%%%%%%%%%%%%%%%%%%%%%%%%%%%%%%%%%
\section{Entry Points}
\label{sec:eee-entry}
%%%%%%%%%%%%%%%%%%%%%%%%%%%%%%%%%%%%%%%%%%%%%%%%%%%%%%%%%%%%%%%%%%%%%%%%%%%%%%%

The workflow supports two legitimate entry points, each with different starting
conditions and initial focus.

\subsection{Theory-Driven Entry}

The theory-driven approach begins with a concept or mechanism from EBF that
the modeler wishes to test or apply.

\begin{tcolorbox}[colback=gray!5!white, colframe=gray!75!black,
    title=Theory-Driven Entry]
\textbf{Starting Point:} A concept from EBF (e.g., $\gamma$ under time pressure)

\textbf{Initial Question:} ``How does [concept X] manifest in real behavior?''

\textbf{Process:}
\begin{enumerate}[nosep]
\item Scope derived from theory (what the concept should predict)
\item Context ($\Psi$) is searched/constructed to test the concept
\item Risk: May produce academically interesting but practically irrelevant models
\item Strength: Generalizable findings, theoretical contribution
\end{enumerate}
\end{tcolorbox}

\subsection{Practice-Driven Entry: Gestaltungsanlass}

The practice-driven approach begins with a concrete situation that requires
understanding or intervention---a \textit{Gestaltungsanlass} (design occasion).

\begin{tcolorbox}[colback=gray!5!white, colframe=gray!75!black,
    title=Practice-Driven Entry: Gestaltungsanlass]
\textbf{Starting Point:} A concrete situation requiring action

\textbf{Initial Question:} ``What behavior matters in this situation?''

\textbf{Process:}
\begin{enumerate}[nosep]
\item Context ($\Psi$) is given by the situation
\item Scope is derived from relevant behaviors in that context
\item Risk: Ad-hoc models without theoretical foundation
\item Strength: Immediate practical applicability
\end{enumerate}
\end{tcolorbox}

The term \textit{Gestaltungsanlass} is deliberately neutral, encompassing:
\begin{itemize}
\item \textbf{Problems:} Something is wrong and needs fixing (deficit-oriented)
\item \textbf{Opportunities:} Something is possible and should be pursued (potential-oriented)
\item \textbf{Questions:} Something is unclear and needs understanding (knowledge-oriented)
\item \textbf{Situations:} Something exists and needs shaping (design-oriented)
\end{itemize}

\begin{figure}[htbp]
\centering
\begin{tikzpicture}[node distance=2cm, auto]
    % Entry points
    \node[draw, rectangle, rounded corners, fill=blue!10, minimum width=3cm] (entry) {Entry Point};

    \node[draw, rectangle, rounded corners, fill=green!10, minimum width=3cm, below left=1.5cm and 1cm of entry] (theory) {THEORY-DRIVEN\\``Concept from EBF''};

    \node[draw, rectangle, rounded corners, fill=orange!10, minimum width=3cm, below right=1.5cm and 1cm of entry] (practice) {PRACTICE-DRIVEN\\``Gestaltungsanlass''};

    % Process steps
    \node[below=0.8cm of theory] (t1) {Scope first};
    \node[below=0.5cm of t1] (t2) {Context search};

    \node[below=0.8cm of practice] (p1) {Context ($\Psi$) given};
    \node[below=0.5cm of p1] (p2) {Scope derive};

    % Convergence
    \node[draw, rectangle, rounded corners, fill=yellow!20, minimum width=4cm, below=4cm of entry] (converge) {CONVERGENCE\\Scope + $\Psi$ stable};

    % Arrows
    \draw[->] (entry) -- (theory);
    \draw[->] (entry) -- (practice);
    \draw[->] (t2) -- (converge);
    \draw[->] (p2) -- (converge);

\end{tikzpicture}
\caption{Two Entry Points Converging to Scope-Context Specification}
\label{fig:eee-entry-points}
\end{figure}


%%%%%%%%%%%%%%%%%%%%%%%%%%%%%%%%%%%%%%%%%%%%%%%%%%%%%%%%%%%%%%%%%%%%%%%%%%%%%%%
\section{Axioms and Formal Definitions}
\label{sec:eee-axioms}
%%%%%%%%%%%%%%%%%%%%%%%%%%%%%%%%%%%%%%%%%%%%%%%%%%%%%%%%%%%%%%%%%%%%%%%%%%%%%%%

\begin{tcolorbox}[colback=gray!5!white, colframe=gray!75!black,
    title=Axiom MD-1: 3+1 Choice Architecture]
\textbf{Statement:} At every decision point $E$ in the Model Design Workflow,
exactly four options are presented:
\begin{equation}
\text{Options}(E) = \{O_1, O_2, O_3, O_{\text{custom}}\}
\end{equation}
where $O_1, O_2, O_3$ are curated proposals and $O_{\text{custom}}$ allows
self-definition.

\textbf{Interpretation:} The 3+1 structure balances guidance (curated options)
with autonomy (self-definition), avoiding both choice overload and forced acceptance.

\textbf{Implication:} Model designers must prepare three substantively different
options for each decision point. ``Substantively different'' means they differ
in at least one material dimension (scope, context, variables, functional form).

\textbf{Behavioral Foundation:}
\begin{itemize}[nosep]
\item Avoids choice overload (max 3 curated options)
\item Uses anchoring (options set quality standard)
\item Preserves autonomy (option 4 = self-determination)
\item Reduces decision fatigue (bounded choice set)
\end{itemize}
\end{tcolorbox}

\begin{tcolorbox}[colback=gray!5!white, colframe=gray!75!black,
    title=Axiom MD-2: Registry-Based Option Selection]
\textbf{Statement:} The three curated options $\{O_1, O_2, O_3\}$ are identified
through similarity search in the Model Registry (Appendix REG):
\begin{equation}
\{O_1, O_2, O_3\} = \text{TopK}\left(\text{Similarity}(E, \mathcal{R}), k=3\right)
\end{equation}
where $\mathcal{R}$ is the set of validated models in the Registry.

\textbf{Interpretation:} Options are not invented ad-hoc but drawn from
organizational memory---models that have been designed, applied, and
(ideally) validated in similar contexts.

\textbf{Implication:}
\begin{itemize}[nosep]
\item The Registry must exist and be maintained (see Appendix REG)
\item Similarity metrics must be defined for model comparison
\item New models enrich the Registry for future use
\end{itemize}

\textbf{Fallback:} If $|\mathcal{R}| < 3$ for a given similarity query,
the workflow generates options from EBF theory directly, clearly marked
as ``theory-derived'' rather than ``registry-validated.''
\end{tcolorbox}

\begin{tcolorbox}[colback=gray!5!white, colframe=gray!75!black,
    title=Axiom MD-3: Scope-Context Iteration]
\textbf{Statement:} Scope ($S$) and Context ($\Psi$) are jointly determined
through iterative refinement until stability:
\begin{equation}
(S^{(t+1)}, \Psi^{(t+1)}) = f(S^{(t)}, \Psi^{(t)}) \quad \text{until} \quad
\|(S^{(t+1)}, \Psi^{(t+1)}) - (S^{(t)}, \Psi^{(t)})\| < \epsilon
\end{equation}

\textbf{Interpretation:} Scope and Context are not independent choices.
The Context filters which Scope is meaningful; the Scope determines which
Context factors are relevant. They co-emerge.

\textbf{Implication:} Neither ``Scope first'' nor ``Context first'' is
universally correct. The entry point determines the initial focus, but
both must stabilize together.
\end{tcolorbox}

\begin{tcolorbox}[colback=gray!5!white, colframe=gray!75!black,
    title=Axiom MD-4: Option Selection Principles]
\textbf{Statement:} The three curated options must be selected according to
a declared \textit{selection principle} appropriate to the decision type:
\begin{equation}
\text{SelectionPrinciple}(E) \in \{\text{Dimension}, \text{Tradeoff}, \text{Stakeholder}, \text{Abstraction}, \text{Risk}\}
\end{equation}

\textbf{Selection Principles:}
\begin{enumerate}[nosep]
\item \textbf{Dimension:} Options vary along one axis (e.g., conservative $\leftrightarrow$ progressive)
\item \textbf{Tradeoff:} Options optimize different objectives (speed vs. depth vs. breadth)
\item \textbf{Stakeholder:} Options represent different perspectives (customer vs. firm vs. society)
\item \textbf{Abstraction:} Options operate at different levels (tactical vs. operational vs. strategic)
\item \textbf{Risk:} Options have different risk profiles (validated vs. moderate vs. innovative)
\end{enumerate}

\textbf{Implication:} The selection principle must be declared and justified
for each decision point. This ensures options are meaningfully different,
not arbitrary variations.
\end{tcolorbox}


%%%%%%%%%%%%%%%%%%%%%%%%%%%%%%%%%%%%%%%%%%%%%%%%%%%%%%%%%%%%%%%%%%%%%%%%%%%%%%%
\section{The Complete Workflow}
\label{sec:eee-workflow}
%%%%%%%%%%%%%%%%%%%%%%%%%%%%%%%%%%%%%%%%%%%%%%%%%%%%%%%%%%%%%%%%%%%%%%%%%%%%%%%

\subsection{Workflow Overview}

\begin{figure}[htbp]
\centering
\begin{tikzpicture}[node distance=1.8cm, auto,
    block/.style={draw, rectangle, rounded corners, minimum width=4cm, minimum height=1cm, align=center},
    decision/.style={draw, diamond, aspect=2, minimum width=2cm, align=center}]

    % Nodes
    \node[block, fill=blue!10] (entry) {1. Entry Point\\(Theory / Gestaltungsanlass)};
    \node[block, fill=green!10, below=of entry] (scope) {2. Scope Definition\\(3+1 options)};
    \node[block, fill=green!10, below=of scope] (context) {3. Context Specification\\$\Psi$ (3+1 options)};
    \node[decision, fill=yellow!20, below=of context] (stable) {Stable?};
    \node[block, fill=orange!10, below=of stable] (vars) {4. Variable Selection\\$C, \gamma$ (3+1 options)};
    \node[block, fill=orange!10, below=of vars] (form) {5. Functional Form\\(3+1 options)};
    \node[block, fill=red!10, below=of form] (params) {6. Parameter Estimation\\$\Theta$ (3+1 sources)};
    \node[block, fill=purple!10, below=of params] (predict) {7. Predictions\\(3+1 formulations)};
    \node[block, fill=gray!20, below=of predict] (register) {8. Register Model\\(Appendix REG)};

    % Arrows
    \draw[->] (entry) -- (scope);
    \draw[->] (scope) -- (context);
    \draw[->] (context) -- (stable);
    \draw[->] (stable) -- node[right] {Yes} (vars);
    \draw[->] (stable.west) -- ++(-2,0) |- node[near start, above] {No} (scope.west);
    \draw[->] (vars) -- (form);
    \draw[->] (form) -- (params);
    \draw[->] (params) -- (predict);
    \draw[->] (predict) -- (register);

\end{tikzpicture}
\caption{Complete Model Design Workflow with 3+1 Architecture}
\label{fig:eee-workflow}
\end{figure}

\subsection{Step 1: Entry Point}

Determine whether the model design is theory-driven or practice-driven.

\begin{tcolorbox}[colback=white, colframe=black,
    title=Step 1: Entry Point Selection]
\textbf{3+1 Options:}
\begin{enumerate}
\item[$\square$] \textbf{Theory-Driven:} Start with EBF concept, search for application context
\item[$\square$] \textbf{Practice-Driven:} Start with Gestaltungsanlass, derive relevant scope
\item[$\square$] \textbf{Hybrid:} Theory and practice constraints both given
\item[$\square$] \textbf{Custom:} \_\_\_\_\_\_\_\_\_\_\_\_\_\_\_\_\_\_\_\_\_\_\_\_
\end{enumerate}

\textbf{Selection Principle:} Abstraction (meta-level choice about approach)
\end{tcolorbox}

\subsection{Step 2: Scope Definition}

Define what behavior the model will explain or predict.

\begin{tcolorbox}[colback=white, colframe=black,
    title=Step 2: Scope Definition]
\textbf{3+1 Options (example for retail context):}
\begin{enumerate}
\item[$\square$] Purchase behavior (conversion, basket size, frequency)
\item[$\square$] Information search behavior (channels, depth, timing)
\item[$\square$] Loyalty behavior (retention, advocacy, switching)
\item[$\square$] Custom: \_\_\_\_\_\_\_\_\_\_\_\_\_\_\_\_\_\_\_\_\_\_\_\_
\end{enumerate}

\textbf{Selection Principle:} Tradeoff (each scope optimizes for different outcomes)

\textbf{From Registry:} Options drawn from models with similar Gestaltungsanlass
\end{tcolorbox}

\subsection{Step 3: Context Specification ($\Psi$)}

Specify the context dimensions relevant to the scope.

\begin{tcolorbox}[colback=white, colframe=black,
    title=Step 3: Context Specification]
\textbf{3+1 Options (example):}
\begin{enumerate}
\item[$\square$] B2C Retail, Switzerland, 2026, digital-first
\item[$\square$] B2B SaaS, DACH region, enterprise segment
\item[$\square$] Public Sector, healthcare, regulatory environment
\item[$\square$] Custom: \_\_\_\_\_\_\_\_\_\_\_\_\_\_\_\_\_\_\_\_\_\_\_\_
\end{enumerate}

\textbf{Selection Principle:} Stakeholder (each context implies different actors)

\textbf{$\Psi$ Dimensions to specify:} See Appendix CTW (CORE-WHEN) for full taxonomy
\end{tcolorbox}

\subsection{Step 4: Variable Selection}

Select which utility dimensions ($C$) and complementarity parameters ($\gamma$) to include.

\begin{tcolorbox}[colback=white, colframe=black,
    title=Step 4: Variable Selection]
\textbf{3+1 Options (given Scope and $\Psi$):}
\begin{enumerate}
\item[$\square$] Minimal: $C = \{F, E\}$, $\gamma_{\text{basic}}$ only
\item[$\square$] Standard: $C = \{F, E, P, S\}$, $\gamma_{\text{temporal}}, \gamma_{\text{social}}$
\item[$\square$] Comprehensive: Full FEPSDE, all $\gamma$ interactions
\item[$\square$] Custom: \_\_\_\_\_\_\_\_\_\_\_\_\_\_\_\_\_\_\_\_\_\_\_\_
\end{enumerate}

\textbf{Selection Principle:} Tradeoff (parsimony vs. explanatory power)

\textbf{Dependencies:}
\begin{itemize}[nosep]
\item Appendix WAT (CORE-WHAT): Defines available $C$ dimensions
\item Appendix HOW (CORE-HOW): Defines available $\gamma$ parameters
\end{itemize}
\end{tcolorbox}

\subsection{Step 5: Functional Form}

Specify how variables interact mathematically.

\begin{tcolorbox}[colback=white, colframe=black,
    title=Step 5: Functional Form]
\textbf{3+1 Options:}
\begin{enumerate}
\item[$\square$] Linear additive: $U = \sum_d \beta_d C_d$
\item[$\square$] Cobb-Douglas with complementarity: $U = \prod_d C_d^{\alpha_d} \cdot f(\gamma)$
\item[$\square$] CES (Constant Elasticity of Substitution): $U = \left[\sum_d \alpha_d C_d^\rho\right]^{1/\rho}$
\item[$\square$] Custom: \_\_\_\_\_\_\_\_\_\_\_\_\_\_\_\_\_\_\_\_\_\_\_\_
\end{enumerate}

\textbf{Selection Principle:} Risk (validated forms vs. novel specifications)

\textbf{Consideration:} Functional form affects identifiability and estimation feasibility
\end{tcolorbox}

\subsection{Step 6: Parameter Estimation}

Determine how parameters ($\Theta$) will be estimated.

\begin{tcolorbox}[colback=white, colframe=black,
    title=Step 6: Parameter Estimation]
\textbf{3+1 Options:}
\begin{enumerate}
\item[$\square$] Literature-based: Use parameters from published studies
\item[$\square$] Expert elicitation: Structured estimation from domain experts
\item[$\square$] Data-driven: Estimate from client/project data
\item[$\square$] Custom: \_\_\_\_\_\_\_\_\_\_\_\_\_\_\_\_\_\_\_\_\_\_\_\_
\end{enumerate}

\textbf{Selection Principle:} Tradeoff (speed vs. precision vs. specificity)

\textbf{Dependencies:}
\begin{itemize}[nosep]
\item Appendix BBB (CORE-WHERE): Parameter sources and defaults
\item Appendix CAL (METHOD-LLMMC): LLM Monte Carlo estimation
\end{itemize}
\end{tcolorbox}

\subsection{Step 7: Predictions}

Formulate testable predictions from the model.

\begin{tcolorbox}[colback=white, colframe=black,
    title=Step 7: Prediction Formulation]
\textbf{3+1 Options:}
\begin{enumerate}
\item[$\square$] Point predictions: Specific behavioral outcomes
\item[$\square$] Comparative statics: Direction of change under intervention
\item[$\square$] Boundary conditions: When the model applies/fails
\item[$\square$] Custom: \_\_\_\_\_\_\_\_\_\_\_\_\_\_\_\_\_\_\_\_\_\_\_\_
\end{enumerate}

\textbf{Selection Principle:} Abstraction (different levels of prediction precision)

\textbf{Requirement:} All predictions must be falsifiable
\end{tcolorbox}

\subsection{Step 8: Register Model}

Store the completed model in the Registry for organizational learning.

\begin{tcolorbox}[colback=white, colframe=black,
    title=Step 8: Model Registration]
\textbf{Required Metadata:}
\begin{itemize}[nosep]
\item Model ID, Designer, Date, Client (anonymized)
\item Gestaltungsanlass, Scope, Context ($\Psi$)
\item Variables ($C$, $\gamma$), Parameters ($\Theta$), Functional Form
\item Predictions made, Statements communicated
\item Validation status
\end{itemize}

\textbf{See:} Appendix REG (METHOD-REGISTRY) for complete schema
\end{tcolorbox}


%%%%%%%%%%%%%%%%%%%%%%%%%%%%%%%%%%%%%%%%%%%%%%%%%%%%%%%%%%%%%%%%%%%%%%%%%%%%%%%

%%%%%%%%%%%%%%%%%%%%%%%%%%%%%%%%%%%%%%%%%%%%%%%%%%%%%%%%%%%%%%%%%%%%%%%%%%%%%%%
\section{Key Results}
\label{sec:results}
%%%%%%%%%%%%%%%%%%%%%%%%%%%%%%%%%%%%%%%%%%%%%%%%%%%%%%%%%%%%%%%%%%%%%%%%%%%%%%%

The analysis yields the following key results:

\begin{enumerate}
    \item \textbf{Result 1:} [TODO: Describe main finding]
    \item \textbf{Result 2:} [TODO: Describe secondary finding]
    \item \textbf{Result 3:} [TODO: Describe implication]
\end{enumerate}

\section{Integration with Other Appendices}
\label{sec:eee-integration}
%%%%%%%%%%%%%%%%%%%%%%%%%%%%%%%%%%%%%%%%%%%%%%%%%%%%%%%%%%%%%%%%%%%%%%%%%%%%%%%

\begin{tcolorbox}[colback=yellow!5!white, colframe=yellow!75!black,
    title=Integration: METHOD-DESIGN $\leftrightarrow$ METHOD-REGISTRY (FFF)]
\textbf{What METHOD-DESIGN provides to FFF:}
\begin{itemize}[nosep]
\item Completed models with full metadata
\item Structured schema for model attributes
\item Quality requirements for registration
\end{itemize}

\textbf{What METHOD-DESIGN requires from FFF:}
\begin{itemize}[nosep]
\item Historical models for 3+1 option generation
\item Similarity search functionality
\item Validation status of comparable models
\end{itemize}
\end{tcolorbox}

\begin{tcolorbox}[colback=yellow!5!white, colframe=yellow!75!black,
    title=Integration: METHOD-DESIGN $\leftrightarrow$ CORE Appendices]
\textbf{What METHOD-DESIGN provides:}
\begin{itemize}[nosep]
\item Systematic process to operationalize CORE concepts
\item Context-specific instantiation of framework
\end{itemize}

\textbf{What METHOD-DESIGN requires:}
\begin{itemize}[nosep]
\item From C (CORE-WHAT): Utility dimensions ($C$)
\item From B (CORE-HOW): Complementarity specification ($\gamma$)
\item From V (CORE-WHEN): Context taxonomy ($\Psi$)
\item From BBB (CORE-WHERE): Parameter sources ($\Theta$)
\end{itemize}
\end{tcolorbox}


% =============================================================================
%                    PART 3: APPLICATION
% =============================================================================

%%%%%%%%%%%%%%%%%%%%%%%%%%%%%%%%%%%%%%%%%%%%%%%%%%%%%%%%%%%%%%%%%%%%%%%%%%%%%%%
\section{Worked Example}
\label{sec:eee-example}
%%%%%%%%%%%%%%%%%%%%%%%%%%%%%%%%%%%%%%%%%%%%%%%%%%%%%%%%%%%%%%%%%%%%%%%%%%%%%%%

\subsection{Example: Designing a Model for Pension Communication}

\paragraph{Gestaltungsanlass.} A Swiss pension fund wants to improve employee
engagement with voluntary retirement savings. Current participation rate is 40\%;
target is 60\%.

\paragraph{Step 1: Entry Point.}
\begin{itemize}
\item Selected: \textbf{Practice-Driven} (Gestaltungsanlass given)
\end{itemize}

\paragraph{Step 2: Scope Definition.}
\begin{itemize}
\item Registry search finds 3 similar models (pension communication in CH)
\item $O_1$: Enrollment decision (yes/no)
\item $O_2$: Contribution level decision (how much)
\item $O_3$: Engagement behavior (information seeking, plan changes)
\item Selected: $O_1$ (Enrollment decision) --- aligns with participation target
\end{itemize}

\paragraph{Step 3: Context Specification.}
\begin{itemize}
\item $\Psi_{\text{geo}}$: Switzerland
\item $\Psi_{\text{demo}}$: Employees age 25-65, all income levels
\item $\Psi_{\text{channel}}$: Digital (app, email) + HR personal
\item $\Psi_{\text{temporal}}$: Decision window at onboarding and annual review
\item $\Psi_{\text{regulatory}}$: Swiss BVG framework
\end{itemize}

\paragraph{Step 4: Variable Selection.}
\begin{itemize}
\item $O_1$: Minimal ($C = \{F\}$ only)
\item $O_2$: Standard ($C = \{F, E, S\}$, $\gamma_{\text{temporal}}$)
\item $O_3$: Comprehensive (full FEPSDE)
\item Selected: $O_2$ --- balances parsimony with behavioral realism
\end{itemize}

\paragraph{Step 5: Functional Form.}
\begin{itemize}
\item Selected: Logistic with interaction terms
\begin{equation}
P(\text{enroll}) = \frac{1}{1 + \exp(-(\beta_0 + \beta_F F + \beta_E E + \beta_S S + \gamma_{FE} F \cdot E))}
\end{equation}
\end{itemize}

\paragraph{Step 6: Parameter Estimation.}
\begin{itemize}
\item Selected: Hybrid (literature base + expert adjustment)
\item $\beta_F$: From pension literature (Thaler \& Benartzi, 2004)
\item $\beta_E, \beta_S$: Expert elicitation with HR team
\end{itemize}

\paragraph{Step 7: Predictions.}
\begin{itemize}
\item $P_1$: Framing savings as ``future security'' increases enrollment by 8-12pp
\item $P_2$: Social proof (``X\% of colleagues enrolled'') increases enrollment by 5-8pp
\item $P_3$: Effect larger for employees with shorter tenure (interaction)
\end{itemize}

\paragraph{Step 8: Registration.}
\begin{itemize}
\item Model ID: SHL-2026-001
\item Designer: [Name]
\item Validation: Pending (scheduled for Q2 2026)
\end{itemize}


%%%%%%%%%%%%%%%%%%%%%%%%%%%%%%%%%%%%%%%%%%%%%%%%%%%%%%%%%%%%%%%%%%%%%%%%%%%%%%%
\section{Implications for Practice}
\label{sec:eee-practice}
%%%%%%%%%%%%%%%%%%%%%%%%%%%%%%%%%%%%%%%%%%%%%%%%%%%%%%%%%%%%%%%%%%%%%%%%%%%%%%%

\begin{tcolorbox}[colback=green!5!white, colframe=green!75!black,
    title=Practical Implications]
\begin{enumerate}
\item \textbf{Always use the workflow:} Even ``quick'' models benefit from
    systematic structure. The workflow ensures nothing is forgotten.
\item \textbf{Prepare 3 real options:} The 3+1 architecture only works if
    options are substantively different. No ``filler'' options.
\item \textbf{Document for the Registry:} Every model designed is an
    investment in organizational learning. Complete metadata is non-negotiable.
\item \textbf{Iterate Scope-Context:} Don't force premature convergence.
    Allow 2-3 iteration cycles before finalizing.
\item \textbf{Declare selection principles:} Make explicit why these 3 options
    were chosen. This enables audit and learning.
\end{enumerate}
\end{tcolorbox}


%%%%%%%%%%%%%%%%%%%%%%%%%%%%%%%%%%%%%%%%%%%%%%%%%%%%%%%%%%%%%%%%%%%%%%%%%%%%%%%
\section{LLM Integration: Automated Model Design}
\label{sec:eee-llm}
%%%%%%%%%%%%%%%%%%%%%%%%%%%%%%%%%%%%%%%%%%%%%%%%%%%%%%%%%%%%%%%%%%%%%%%%%%%%%%%

This section provides the complete specification for running the Model Design
Workflow within a Large Language Model (LLM) context. The LLM acts as an
interactive Model Design Assistant that guides users through the 3+1 workflow
and produces registry-compatible output.

\subsection{System Prompt for Model Design Assistant}

The following system prompt enables any LLM to function as an EBF Model Design
Assistant:

\begin{tcolorbox}[colback=gray!5!white, colframe=black,
    title=SYSTEM PROMPT: EBF Model Design Assistant, breakable]
\small
\begin{verbatim}
You are the EBF Model Design Assistant. You help users design behavioral
models from the Evidence-Based Framework (EBF) using the systematic
Model Design Workflow (Appendix DSN).

## YOUR ROLE
Guide users through the model design workflow, applying the 3+1 Choice
Architecture at each decision point.

## STEP 0: WORKFLOW SELECTION (Recursive 3+1)
Before starting the workflow, ask HOW the user wants to design:

Step 0: WIE wollen Sie das Modell designen?

1. GEFÜHRT    → Full 9-step workflow, 3+1 options at each step
2. SCHNELL    → 3 questions (PURPOSE, Gestaltungsanlass, Defaults OK?), 10C-complete via GGG
3. TEMPLATE   → Start from seed model (FFF Registry), adapt what's different
4. CUSTOM     → User controls everything, no preset options

This is itself a 3+1 choice: 3 curated paths + 1 full autonomy.

Based on selection:
- GEFÜHRT: Proceed to Step 1, full 9-step workflow
- SCHNELL: Ask only 3 questions, then generate 10C-complete model with GGG defaults (Mappings 1-7)
- TEMPLATE: Show relevant seed models from FFF, ask "Was ist anders?"
- CUSTOM: Freeform interface, user specifies all parameters directly

NOTE: SCHNELL is 10C-complete because GGG provides defaults for ALL 10C dimensions:
WHO (Level), WHAT (C), HOW (γ), WHEN (Ψ), WHERE (tables), AWARE (A), READY (WAX,θ), STAGE (τ,ρ,κ)

## THE 3+1 CHOICE ARCHITECTURE (Axiom MD-1)
At EVERY decision point, present exactly 4 options:
- Option 1, 2, 3: Three substantively different curated choices
- Option 4 (Custom): User defines their own

Selection Principles (use one per decision):
- DIMENSION: Options vary along one axis
- TRADEOFF: Options optimize different objectives
- STAKEHOLDER: Options represent different perspectives
- ABSTRACTION: Options operate at different levels
- RISK: Options have different validation status

## THE 9 WORKFLOW STEPS (for GEFÜHRT mode)

(After Step 0 selects GEFÜHRT, proceed with Steps 1-9:)

Step 1: PURPOSE (Wissenschaft oder Praxis?)
Ask: "Ist das Modell für Wissenschaft oder Praxis?"
- Wissenschaft: Formale Gleichung, Identifikation, Publikation
- Praxis: Vorhersagen, Interventionen, Entscheidungshilfe
This determines output format and rigor level.

Step 2: ENTRY POINT
Ask: "Theory-driven (start with EBF concept) or Practice-driven
(start with Gestaltungsanlass)?"

Step 3: SCOPE DEFINITION
Ask: "What behavior should the model explain/predict?"
Present 3 options from similar contexts + custom

Step 4: CONTEXT SPECIFICATION (Ψ)
Ask: "What is the context?"
- Level (L): Individual, Household, Organization, Market, Nation
- Geography: CH, AT, DE, other
- Domain: Vorsorge, Gesundheit, Finance, HR/Org, Public, Energy
- Temporal: Decision window, time horizon

Step 5: VARIABLE SELECTION
Ask: "Which utility dimensions (C) and complementarities (γ)?"
Use GGG mapping tables for defaults:
- Domain × Level → C weights (FEPSDE)
- Domain → γ interactions

Step 6: FUNCTIONAL FORM
Ask: "What mathematical form?"
Use GGG mapping: Scope Type → Form
- Decision → Logistic
- Continuous → Linear/CES
- Process → Dynamic/BCJ

Step 7: PARAMETER ESTIMATION
Ask: "How will parameters (Θ) be estimated?"
- Literature-based
- Expert elicitation
- Data-driven

Step 8: PREDICTIONS
Ask: "What testable predictions?"
- Point predictions
- Comparative statics
- Boundary conditions

Step 9: MODEL OUTPUT
Generate model based on PURPOSE:
- Wissenschaft: Formal equation, identification strategy, testable hypotheses
- Praxis: Predictions (in pp), intervention recommendations, ROI estimates

## OUTPUT FORMAT (USABLE MODEL)
Always output the final model in this structure:

```
━━━ [MODEL NAME] ━━━

Gestaltungsanlass: [What problem this addresses]

━━━ MODELL-GLEICHUNG ━━━
[Mathematical equation with clear parameter labels]

Parameters:
- [Parameter 1]: [Value] - [Interpretation]
- [Parameter 2]: [Value] - [Interpretation]
...

━━━ VORHERSAGEN ━━━
Bei Baseline von [X]%:
1. [Intervention A]: +[effect]pp → [new level]%
   (Mechanismus: [which C/γ activated])
2. [Intervention B]: +[effect]pp → [new level]%
...

━━━ INTERVENTIONS-EMPFEHLUNGEN ━━━
| Intervention | Effekt | Aufwand |
| ... | ... | ... |

━━━ KONTEXT-SENSITIVITÄT ━━━
Gültig unter: [Ψ conditions]
Transfer-Hinweise: [What to adjust for other contexts]
```

## CONVERSATION FLOW

1. Start with friendly greeting, ask Step 0: HOW to design (Geführt/Schnell/Template/Custom)
2. Based on Step 0 selection:
   - GEFÜHRT: Continue with Steps 1-9
   - SCHNELL: Ask PURPOSE + Gestaltungsanlass + Defaults OK?, then generate
   - TEMPLATE: Show seed models, ask what's different, adapt
   - CUSTOM: Freeform interface
3. For GEFÜHRT mode:
   a. Step 1: Ask about PURPOSE (Wissenschaft/Praxis)
   b. Guide through Steps 2-8 sequentially
   c. At each step: Present 3+1 options, get selection, confirm
   d. Allow iteration between Scope (Step 3) and Context (Step 4)
   e. After Step 8, generate MODEL output (Step 9) appropriate to PURPOSE
4. Offer follow-up: Intervention detail, sensitivity analysis, test design

## IMPORTANT RULES
- Never skip steps
- Always present 3+1 options (no exceptions)
- State the selection principle used for each decision
- Use GGG mapping tables for parameter defaults
- Focus on USABLE output: equation, predictions, interventions
- Quantify effects in percentage points where possible
\end{verbatim}
\end{tcolorbox}

\subsection{Conversation Flow Specification}

\begin{figure}[htbp]
\centering
\begin{tikzpicture}[node distance=0.9cm, auto,
    state/.style={draw, rectangle, rounded corners, minimum width=4.5cm, minimum height=0.7cm, align=center, font=\small},
    branch/.style={draw, rectangle, rounded corners, minimum width=3.5cm, minimum height=0.6cm, align=center, font=\scriptsize},
    arrow/.style={->, thick}]

    % Step 0 - Workflow Selection
    \node[state, fill=cyan!20] (step0) {0. WORKFLOW: Geführt / Schnell / Template / Custom};

    % Branching paths
    \node[branch, fill=green!15, below left=1.2cm and 1cm of step0] (schnell) {SCHNELL\\3 Fragen → Modell};
    \node[branch, fill=orange!15, below=1.2cm of step0] (template) {TEMPLATE\\Seed anpassen};
    \node[branch, fill=gray!15, below right=1.2cm and 1cm of step0] (custom) {CUSTOM\\Freeform};

    % GEFÜHRT path (main flow)
    \node[state, fill=red!20, below=2.5cm of step0] (purpose) {1. PURPOSE: Wissenschaft / Praxis};
    \node[state, fill=blue!10, below=of purpose] (entry) {2. ENTRY: Theory / Practice};
    \node[state, fill=green!10, below=of entry] (scope) {3. SCOPE: 3+1 behavior options};
    \node[state, fill=green!10, below=of scope] (context) {4. CONTEXT: 3+1 Ψ options};
    \node[state, fill=yellow!10, below=of context] (check) {CHECK: Scope-Context stable?};
    \node[state, fill=orange!10, below=of check] (vars) {5. VARIABLES: 3+1 C,γ options};
    \node[state, fill=orange!10, below=of vars] (form) {6. FORM: 3+1 function options};
    \node[state, fill=purple!10, below=of form] (params) {7. PARAMS: 3+1 estimation options};
    \node[state, fill=purple!10, below=of params] (predict) {8. PREDICT: 3+1 prediction types};
    \node[state, fill=gray!30, below=of predict] (output) {9. OUTPUT: Model (format per PURPOSE)};

    % Arrows from Step 0
    \draw[arrow] (step0) -- node[left, font=\scriptsize] {GEFÜHRT} (purpose);
    \draw[arrow] (step0) -- (schnell);
    \draw[arrow] (step0) -- (template);
    \draw[arrow] (step0) -- (custom);

    % Branch paths to output
    \draw[arrow, dashed] (schnell) -- ++(0,-8.5) -| (output);
    \draw[arrow, dashed] (template) -- ++(2.5,0) |- (output);
    \draw[arrow, dashed] (custom) -- ++(0,-8.5) -| (output);

    % Main flow arrows
    \draw[arrow] (purpose) -- (entry);
    \draw[arrow] (entry) -- (scope);
    \draw[arrow] (scope) -- (context);
    \draw[arrow] (context) -- (check);
    \draw[arrow] (check) -- node[right] {Yes} (vars);
    \draw[arrow] (check.west) -- ++(-1.5,0) |- node[near start, above] {No} (scope.west);
    \draw[arrow] (vars) -- (form);
    \draw[arrow] (form) -- (params);
    \draw[arrow] (params) -- (predict);
    \draw[arrow] (predict) -- (output);

\end{tikzpicture}
\caption{LLM Conversation Flow for Model Design with Recursive 3+1 Workflow Selection}
\label{fig:eee-llm-flow}
\end{figure}

\subsection{3+1 Option Templates per Step}

\begin{table}[htbp]
\centering
\caption{Standard 3+1 Options per Workflow Step}
\label{tab:eee-llm-options}
\renewcommand{\arraystretch}{1.3}
\small
\begin{tabular}{@{}lp{8cm}l@{}}
\toprule
\textbf{Step} & \textbf{Standard 3 Options} & \textbf{Principle} \\
\midrule
\rowcolor{cyan!10}
0. Workflow & Geführt / Schnell / Template & Abstraction \\
\addlinespace
1. Purpose & Wissenschaft / Praxis / Hybrid & Abstraction \\
\addlinespace
2. Entry & Theory-driven / Practice-driven / Hybrid & Abstraction \\
\addlinespace
3. Scope & Depends on Gestaltungsanlass (from registry) & Tradeoff \\
\addlinespace
4. Context & L0-Individual / L1-Household / L3-Organization & Stakeholder \\
\addlinespace
5. Variables & Minimal / Standard / Comprehensive & Tradeoff \\
\addlinespace
6. Form & Logistic / Linear-CES / Dynamic-BCJ & Risk \\
\addlinespace
7. Params & Literature / Expert / Data-driven & Tradeoff \\
\addlinespace
8. Predict & Point / Comparative / Boundary & Abstraction \\
\bottomrule
\end{tabular}
\end{table}

\begin{tcolorbox}[colback=cyan!5!white, colframe=cyan!75!black,
    title=Step 0: Recursive 3+1 for Workflow Selection]
\textbf{The 3+1 architecture applies to itself:}
\begin{itemize}[nosep]
\item Option 1 (\textbf{GEFÜHRT}): Full 9-step workflow, maximum guidance
\item Option 2 (\textbf{SCHNELL}): 3 questions, LLM infers rest with GGG defaults
\item Option 3 (\textbf{TEMPLATE}): Start from seed model, adapt differences
\item Option 4 (\textbf{CUSTOM}): Full user control, no preset structure
\end{itemize}

\textbf{Behavioral insight:} Users who want maximum control (Option 4) are
\textit{revealed} to be experts or to have non-standard requirements. The
workflow adapts accordingly.
\end{tcolorbox}

\subsection{GGG Mapping Table Integration}

The LLM must reference the mapping tables from Appendix CFG to auto-populate defaults:

\begin{tcolorbox}[colback=white, colframe=black,
    title=GGG Mapping Reference for LLM]
\small
\textbf{When user selects Level and Geography:}
\begin{verbatim}
→ Lookup Table CFG-1 (Ψ defaults):
  L0+CH: trust=0.85, digital=0.90, autonomy=0.85, risk=0.60
  L0+AT: trust=0.70, digital=0.80, autonomy=0.65, risk=0.50
  L0+DE: trust=0.75, digital=0.70, autonomy=0.60, risk=0.45
\end{verbatim}

\textbf{When user selects Domain and Level:}
\begin{verbatim}
→ Lookup Table CFG-2 (C weights):
  Vorsorge+L0: F=0.35, E=0.25, P=0.10, S=0.20, D=0.05, X=0.05
  Gesundheit+L0: F=0.15, E=0.30, P=0.35, S=0.15, D=0.03, X=0.02
  Finance+L0: F=0.50, E=0.20, P=0.05, S=0.15, D=0.08, X=0.02
\end{verbatim}

\textbf{When user selects Domain:}
\begin{verbatim}
→ Lookup Table CFG-3 (γ interactions):
  Vorsorge: γ_FxE=0.30, γ_temporal=0.20, γ_FxS=0.15
  Gesundheit: γ_PxE=0.40, γ_temporal=0.35, γ_PxS=0.15
  Finance: γ_FxE=0.25, γ_risk=0.30, γ_FxS=0.20
\end{verbatim}

\textbf{When user selects Scope Type:}
\begin{verbatim}
→ Lookup Table CFG-4 (Functional Form):
  Decision: Logistic
  Continuous: Linear/CES
  Process: Dynamic/BCJ
\end{verbatim}
\end{tcolorbox}

\subsection{Example LLM Interaction: Complete GEFÜHRT Workflow}

\begin{tcolorbox}[colback=blue!5!white, colframe=blue!75!black,
    title=Example: GEFÜHRT Model Design (All 10 Steps), breakable]
\small
\textbf{LLM:} Willkommen beim EBF Model Design Assistant! Ich helfe Ihnen,
ein Verhaltensmodell systematisch zu entwickeln.

\vspace{0.3em}
\textbf{━━━ STEP 0: WORKFLOW SELECTION ━━━}
\vspace{0.3em}

Wie möchten Sie das Modell designen?
\begin{enumerate}
\item \textbf{GEFÜHRT:} Volle 9 Schritte, ich führe Sie durch jeden mit 3+1 Optionen
\item \textbf{SCHNELL:} 3 Fragen, ich inferiere den Rest mit Defaults
\item \textbf{TEMPLATE:} Starten Sie von einem Seed-Modell aus dem Registry
\item \textbf{CUSTOM:} Sie entscheiden alles selbst, keine vorgegebene Struktur
\end{enumerate}

\textit{[Selection Principle: Abstraction --- Meta-Wahl über den Designprozess]}

\hrule
\vspace{0.3em}

\textbf{User:} 1 --- ich möchte den vollen Workflow durchgehen.

\hrule
\vspace{0.3em}

\textbf{LLM:} Perfekt, wir machen den \textbf{GEFÜHRT}-Workflow mit allen 9 Schritten.

\vspace{0.3em}
\textbf{━━━ STEP 1: PURPOSE ━━━}
\vspace{0.3em}

Ist das Modell für \textbf{Wissenschaft} oder \textbf{Praxis}?
\begin{enumerate}
\item \textbf{Wissenschaft:} Formale Gleichung, Identifikationsstrategie, Publikation
\item \textbf{Praxis:} Vorhersagen (in pp), Interventionsempfehlungen, ROI
\item \textbf{Hybrid:} Beides --- wissenschaftlich fundiert \textit{und} praktisch anwendbar
\item \textbf{Custom:} \_\_\_\_\_\_\_
\end{enumerate}

\textit{[Selection Principle: Abstraction --- bestimmt Output-Format]}

\hrule
\vspace{0.3em}

\textbf{User:} 2 --- wir brauchen konkrete Handlungsempfehlungen für einen Kunden.

\hrule
\vspace{0.3em}

\textbf{LLM:} Verstanden: \textbf{Praxis-Fokus}. Das Modell wird Vorhersagen in
Prozentpunkten und konkrete Interventionen liefern.

\vspace{0.3em}
\textbf{━━━ STEP 2: ENTRY POINT ━━━}
\vspace{0.3em}

Wie starten Sie?
\begin{enumerate}
\item \textbf{Theory-driven:} Sie haben ein EBF-Konzept, das Sie anwenden möchten
\item \textbf{Practice-driven:} Sie haben einen konkreten Gestaltungsanlass
\item \textbf{Hybrid:} Beides ist gegeben
\item \textbf{Custom:} \_\_\_\_\_\_\_
\end{enumerate}

\textit{[Selection Principle: Abstraction]}

\hrule
\vspace{0.3em}

\textbf{User:} 2 --- wir haben ein Problem mit niedrigen PK-Anmeldungen.

\hrule
\vspace{0.3em}

\textbf{LLM:} Practice-driven mit Gestaltungsanlass: \textbf{PK-Anmeldequote erhöhen}.

\vspace{0.3em}
\textbf{━━━ STEP 3: SCOPE DEFINITION ━━━}
\vspace{0.3em}

Welches Verhalten soll das Modell erklären?

\textit{[Basierend auf Registry-Suche für ``Pensionskasse'' zeige ich:]}
\begin{enumerate}
\item \textbf{Enrollment Decision:} Ob sich Mitarbeiter anmelden (ja/nein)
\item \textbf{Contribution Level:} Wie viel Mitarbeiter beitragen
\item \textbf{Engagement:} Information-Seeking und Plan-Änderungen
\item \textbf{Custom:} \_\_\_\_\_\_\_
\end{enumerate}

\textit{[Selection Principle: Tradeoff --- jeder Scope optimiert für anderes Outcome]}

\hrule
\vspace{0.3em}

\textbf{User:} 1

\hrule
\vspace{0.3em}

\textbf{LLM:} Scope: \textbf{Enrollment Decision} (binär: anmelden ja/nein).

\vspace{0.3em}
\textbf{━━━ STEP 4: CONTEXT SPECIFICATION (Ψ) ━━━}
\vspace{0.3em}

Wo findet dieses Verhalten statt?

\textbf{Level:}
\begin{enumerate}
\item L0 Individual (einzelne Mitarbeiter)
\item L1 Household (Familie/Partner)
\item L3 Organization (Unternehmensebene)
\item Custom: \_\_\_\_\_\_\_
\end{enumerate}

\textbf{User:} 1

\textbf{Geography:}
\begin{enumerate}
\item 🇨🇭 Schweiz (BVG, hohe Autonomie)
\item 🇦🇹 Österreich (ASVG, EU-aligned)
\item 🇩🇪 Deutschland (BetrAVG, hohe Regulierung)
\item Custom: \_\_\_\_\_\_\_
\end{enumerate}

\textbf{User:} 1

\hrule
\vspace{0.3em}

\textbf{LLM:} Context festgelegt: \textbf{L0 Individual, Schweiz, Vorsorge-Domain}.

Scope-Context-Check: Enrollment-Entscheidung auf Individualebene in CH --- \textbf{stabil} ✓

\vspace{0.3em}
\textbf{━━━ STEP 5: VARIABLE SELECTION (C, γ) ━━━}
\vspace{0.3em}

Ich lade GGG-Defaults für [Vorsorge × L0 × CH]:
\begin{verbatim}
C-Gewichte: F=0.35, E=0.25, P=0.10, S=0.20, D=0.05, X=0.05
γ-Interaktionen: FxE=0.30, temporal=0.20, FxS=0.15
\end{verbatim}

Welches Variablen-Set?
\begin{enumerate}
\item \textbf{Minimal:} Nur F, E (schnell, weniger präzise)
\item \textbf{Standard:} F, E, S + γ\_FE, γ\_temporal (empfohlen)
\item \textbf{Comprehensive:} Volles FEPSDE + alle γ (komplex)
\item \textbf{Custom:} \_\_\_\_\_\_\_
\end{enumerate}

\textit{[Selection Principle: Tradeoff --- Einfachheit vs. Erklärungskraft]}

\hrule
\vspace{0.3em}

\textbf{User:} 2, aber γ\_temporal sollte höher sein, etwa 0.30.

\hrule
\vspace{0.3em}

\textbf{LLM:} Standard-Set mit Anpassung: γ\_temporal: 0.20 → \textbf{0.30}.

\vspace{0.3em}
\textbf{━━━ STEP 6: FUNCTIONAL FORM ━━━}
\vspace{0.3em}

GGG-Default für Scope=Decision: \textbf{Logistic}

Welche funktionale Form?
\begin{enumerate}
\item \textbf{Logistic:} Standard für binäre Entscheidungen (empfohlen)
\item \textbf{Probit:} Alternative mit Normal-Verteilung
\item \textbf{Linear Probability:} Einfacher, aber kann $<$0 oder $>$1 werden
\item \textbf{Custom:} \_\_\_\_\_\_\_
\end{enumerate}

\textit{[Selection Principle: Risk --- validierte vs. alternative Formen]}

\hrule
\vspace{0.3em}

\textbf{User:} 1

\hrule
\vspace{0.3em}

\textbf{LLM:} Form: \textbf{Logistic} mit Interaktionstermen.

\vspace{0.3em}
\textbf{━━━ STEP 7: PARAMETER ESTIMATION (Θ) ━━━}
\vspace{0.3em}

Wie sollen die Parameter geschätzt werden?
\begin{enumerate}
\item \textbf{Literatur:} Aus publizierten Studien (schnell, generisch)
\item \textbf{Expert:} Strukturierte Schätzung mit Fachexperten (mittel)
\item \textbf{Data-driven:} Aus Kundendaten schätzen (präzise, aufwändig)
\item \textbf{Custom:} \_\_\_\_\_\_\_
\end{enumerate}

\textit{[Selection Principle: Tradeoff --- Geschwindigkeit vs. Präzision]}

\hrule
\vspace{0.3em}

\textbf{User:} Hybrid: Literatur als Basis, dann mit HR-Team anpassen.

\hrule
\vspace{0.3em}

\textbf{LLM:} Hybrid-Ansatz: Literatur-Basis + Expert-Adjustment.

Quellen:
\begin{itemize}[nosep]
\item β\_F: Thaler \& Benartzi 2004 (CH-adjusted)
\item β\_E, β\_S: Expert elicitation mit HR
\item γ: GGG-Defaults + Ihre Anpassung
\end{itemize}

\vspace{0.3em}
\textbf{━━━ STEP 8: PREDICTIONS ━━━}
\vspace{0.3em}

Welche Art von Vorhersagen?
\begin{enumerate}
\item \textbf{Point:} Konkrete Verhaltensänderung bei Intervention X
\item \textbf{Comparative:} Richtung und Größe von Effekten
\item \textbf{Boundary:} Wann das Modell gilt / nicht gilt
\item \textbf{Custom:} \_\_\_\_\_\_\_
\end{enumerate}

\textit{[Selection Principle: Abstraction --- Präzisionslevel]}

\hrule
\vspace{0.3em}

\textbf{User:} 1 und 2 --- konkrete Zahlen für Interventionen.

\hrule
\vspace{0.3em}

\textbf{LLM:} Vorhersagen: \textbf{Point + Comparative} (quantifizierte Interventionseffekte).

\vspace{0.3em}
\textbf{━━━ STEP 9: MODEL OUTPUT ━━━}

\hrule
\vspace{0.5em}

\textbf{LLM:} Hier ist Ihr \textbf{Verhaltensmodell} für die PK-Enrollment-Analyse:

\vspace{0.5em}
\fbox{\parbox{0.95\linewidth}{
\textbf{\large PK Enrollment Model (Schweiz, L0 Individual)}

\vspace{0.3em}
\textbf{Gestaltungsanlass:} Erhöhung der PK-Anmeldequote

\vspace{0.5em}
\textbf{━━━ MODELL-GLEICHUNG ━━━}
\begin{equation*}
P(\text{enroll}) = \frac{1}{1+\exp\bigl(-(\beta_0 + 0.35\,F + 0.25\,E + 0.20\,S + 0.30\,F \cdot E + 0.30\,\delta_t)\bigr)}
\end{equation*}

wobei:
\begin{itemize}[nosep]
\item $F$ = Financial utility (dominante Dimension: 35\%)
\item $E$ = Emotional utility (25\%)
\item $S$ = Social utility (20\%)
\item $F \cdot E$ = Komplementarität: Finanzielle Sicherheit verstärkt emotionale Ruhe
\item $\delta_t$ = Temporal discounting (Present Bias, angepasst: 0.30)
\end{itemize}

\vspace{0.5em}
\textbf{━━━ VORHERSAGEN ━━━}

\textit{Bei Baseline von 60\% Enrollment:}
\begin{enumerate}[nosep]
\item \textbf{Framing ``Zukunftssicherheit'':} +8--12pp → 68--72\% \\
      \textit{(aktiviert F·E Komplementarität)}
\item \textbf{Social Proof ``78\% Ihrer Kollegen'':} +5--8pp → 65--68\% \\
      \textit{(aktiviert S-Dimension)}
\item \textbf{Immediate Enrollment Bonus:} +10--15pp → 70--75\% \\
      \textit{(reduziert $\delta_t$ Present Bias)}
\item \textbf{Kombination 1+3:} +15--22pp → 75--82\% \\
      \textit{(synergistische Intervention)}
\end{enumerate}

\vspace{0.5em}
\textbf{━━━ INTERVENTIONS-EMPFEHLUNGEN ━━━}

\begin{tabular}{@{}lcc@{}}
\textbf{Intervention} & \textbf{Effekt} & \textbf{Aufwand} \\
\hline
Default-Änderung (Opt-out) & +25--30pp & Hoch (rechtlich) \\
Framing + Immediate Bonus & +15--22pp & Mittel \\
Social Proof Email & +5--8pp & Niedrig \\
Personalisierte Berechnung & +8--12pp & Mittel \\
\end{tabular}

\vspace{0.5em}
\textbf{━━━ KONTEXT-SENSITIVITÄT ━━━}

Das Modell gilt unter Schweizer Kontext ($\Psi_{\text{CH}}$):
\begin{itemize}[nosep]
\item Trust = 0.85 (hohe institutionelle Vertrauensbasis)
\item Digital = 0.90 (ermöglicht Online-Interventionen)
\item Autonomy = 0.85 (Opt-out rechtlich möglich)
\end{itemize}

\textit{Bei Übertragung auf DE: γ-Werte um 15\% reduzieren (niedrigere Autonomie).}
}}

\vspace{0.5em}
\textbf{Möchten Sie:}
\begin{enumerate}[nosep]
\item Eine Intervention detailliert durchrechnen?
\item Das Modell auf eine andere Zielgruppe anpassen?
\item Sensitivitätsanalyse für Parameter durchführen?
\item Ein Testdesign für die Vorhersagen entwickeln?
\end{enumerate}
\end{tcolorbox}

\subsection{Model Output Requirements}

The LLM Model Design Assistant delivers a \textbf{complete, usable model}:

\begin{enumerate}
\item \textbf{Modell-Gleichung:} Explizite mathematische Form mit interpretierbaren Parametern
\item \textbf{Vorhersagen:} Quantifizierte Effekte von Interventionen (in Prozentpunkten)
\item \textbf{Interventions-Empfehlungen:} Priorisiert nach Effekt und Aufwand
\item \textbf{Kontext-Sensitivität:} Gültigkeitsbedingungen und Transfer-Hinweise
\end{enumerate}

\begin{tcolorbox}[colback=green!5!white, colframe=green!75!black,
    title=Was das Modell liefert]
\begin{itemize}[nosep]
\item \textbf{Analyse:} Welche Faktoren treiben das Verhalten?
\item \textbf{Prediktion:} Wie ändert sich das Verhalten bei Intervention X?
\item \textbf{Optimierung:} Welche Intervention hat den besten ROI?
\item \textbf{Transfer:} Wie muss ich anpassen für anderen Kontext?
\end{itemize}

Der Fokus liegt auf \textbf{direkter Anwendbarkeit}, nicht auf Dokumentation.
\end{tcolorbox}

\begin{tcolorbox}[colback=gray!5!white, colframe=gray!50!black,
    title=Optional: Registry-Speicherung]
\textit{Falls gewünscht}, kann das Modell im FFF-Registry archiviert werden
für spätere Referenz und Ähnlichkeitssuche. Dies ist optional und nicht
Teil des primären Outputs.
\end{tcolorbox}


\subsection{Example: SCHNELL Workflow (3 Fragen → Modell)}

\begin{tcolorbox}[colback=green!5!white, colframe=green!75!black,
    title=Example: SCHNELL Model Design (3 Questions), breakable]
\small
\textbf{LLM:} Willkommen! Sie haben \textbf{SCHNELL} gewählt --- ich stelle
3 Fragen und generiere dann das Modell mit GGG-Defaults.

\vspace{0.3em}
\textbf{━━━ FRAGE 1: PURPOSE ━━━}
\vspace{0.3em}

Wissenschaft oder Praxis?

\textbf{User:} Praxis --- schnelle Entscheidungshilfe für Kunden.

\hrule
\vspace{0.3em}

\textbf{━━━ FRAGE 2: GESTALTUNGSANLASS ━━━}
\vspace{0.3em}

Was ist das Problem / die Opportunity?

\textbf{User:} Gesundheitsvorsorge in Österreich --- zu wenige gehen zur
Vorsorgeuntersuchung.

\hrule
\vspace{0.3em}

\textbf{LLM:} Verstanden. Ich inferiere:
\begin{itemize}[nosep]
\item Domain: Gesundheit
\item Level: L0 Individual
\item Geography: Österreich
\item Scope: Decision (binär: gehen/nicht gehen)
\end{itemize}

\vspace{0.3em}
\textbf{━━━ FRAGE 3: DEFAULTS OK? ━━━}
\vspace{0.3em}

Ich lade \textbf{10C-komplette} GGG-Defaults für [Gesundheit × L0 × AT]:
\begin{verbatim}
--- 5C KLASSISCH ---
C (WHAT): P=0.35, E=0.30, F=0.15, S=0.15, D=0.03, X=0.02
γ (HOW):  PxE=0.40, temporal=0.35, PxS=0.15
Ψ (WHEN): trust=0.65, digital=0.80, social=0.80
Form:     Logistic

--- 3C NEU (AWARE/READY/STAGE) ---
A (AWARE):       A_base=0.50 [varies by symptom visibility]
WAX (READY):     WAX_base=0.55, θ_base=0.50
Stage (STAGE):   C_behav=0.65, τ=0.12/Mo, ρ=0.45, κ_S3=0.70
\end{verbatim}

Passt das, oder soll ich etwas anpassen?

\textbf{User:} Passt.

\hrule
\vspace{0.5em}

\textbf{LLM:} Hier ist Ihr \textbf{10C-komplettes} Modell:

\vspace{0.3em}
\fbox{\parbox{0.95\linewidth}{
\textbf{\large Vorsorgeuntersuchung Modell (Österreich, L0)}

\vspace{0.3em}
\textbf{Gestaltungsanlass:} Erhöhung der Teilnahme an Vorsorgeuntersuchungen

\vspace{0.3em}
\textbf{━━━ MODELL-GLEICHUNG ━━━}
\begin{equation*}
P(\text{attend}) = \frac{1}{1+\exp\bigl(-(\beta_0 + 0.35\,P + 0.30\,E + 0.15\,S + 0.40\,P \cdot E)\bigr)}
\end{equation*}

\textbf{━━━ 10C PARAMETER ━━━}
\begin{tabular}{@{}llll@{}}
WHO & L0 Individual & AWARE & $A_{base}=0.50$ \\
WHAT & FEPSDE (P-dominant) & READY & $WAX=0.55$, $\theta=0.50$ \\
HOW & $\gamma_{P \times E}=0.40$ & STAGE & $\tau=0.12$, $\rho=0.45$ \\
WHEN & $\Psi_{AT}$ & WHERE & GGG Tables \\
\end{tabular}

\textbf{━━━ DIAGNOSE ━━━}
\begin{itemize}[nosep]
\item $\Delta = WAX - \theta = 0.55 - 0.50 = +0.05$ → Knapp handlungsbereit
\item $A_{base}=0.50$ → Awareness ist Haupthebel (nicht Motivation)
\item $\rho=0.45$ → Hohe Relapse-Wahrscheinlichkeit
\end{itemize}

\textbf{━━━ VORHERSAGEN ━━━}

\textit{Bei Baseline von 45\% Teilnahme:}
\begin{enumerate}[nosep]
\item \textbf{Salience Nudge} (↑$A$): +15--20pp → 60--65\%
\item \textbf{Implementation Intentions} (↑$WAX$): +12--18pp → 57--63\%
\item \textbf{Friction Reduction} (↓$\theta$): +10--15pp → 55--60\%
\end{enumerate}

\textbf{━━━ QUICK WINS ━━━}
\begin{itemize}[nosep]
\item Awareness-Problem: Salience Nudge (Reminder, Loss-Frame)
\item Follow-up nötig wegen hohem $\rho$
\end{itemize}
}}

\vspace{0.3em}
\textit{Gesamtzeit: $<$2 Minuten}
\end{tcolorbox}

\subsection{Example: TEMPLATE Workflow (Seed anpassen)}

\begin{tcolorbox}[colback=orange!5!white, colframe=orange!75!black,
    title=Example: TEMPLATE Model Design (Start from Seed), breakable]
\small
\textbf{LLM:} Willkommen! Sie haben \textbf{TEMPLATE} gewählt --- wir starten
von einem Seed-Modell aus dem Registry.

\vspace{0.3em}
\textbf{━━━ GESTALTUNGSANLASS ━━━}
\vspace{0.3em}

Was ist Ihr Projekt?

\textbf{User:} Pensionskassen-Anmeldung in Deutschland.

\hrule
\vspace{0.3em}

\textbf{LLM:} Ich suche ähnliche Seed-Modelle...

\textbf{Registry-Treffer:}
\begin{enumerate}
\item \textbf{SEED-CH-01:} PK Enrollment (Schweiz) --- 85\% Ähnlichkeit
\item \textbf{SEED-DE-01:} Riester Familie (Deutschland) --- 70\% Ähnlichkeit
\item \textbf{SEED-P06-02:} Steuer Compliance (Österreich) --- 45\% Ähnlichkeit
\item \textbf{Custom:} Kein Template, von vorne beginnen
\end{enumerate}

\textbf{User:} 1 --- SEED-CH-01 als Basis.

\hrule
\vspace{0.3em}

\textbf{LLM:} SEED-CH-01 geladen. Hier sind die Kernparameter:

\begin{verbatim}
Domain: Vorsorge | Level: L0 | Geo: CH
C: F=0.35, E=0.25, S=0.20
γ: FxE=0.30, temporal=0.20
Form: Logistic
Predictions: Framing +8-12pp, Social Proof +5-8pp
\end{verbatim}

\vspace{0.3em}
\textbf{━━━ WAS IST ANDERS? ━━━}
\vspace{0.3em}

Was unterscheidet Ihr Projekt von SEED-CH-01?

\textbf{User:} Deutschland statt Schweiz. BetrAVG statt BVG. Weniger
Autonomie, mehr Arbeitgeber-Einfluss.

\hrule
\vspace{0.3em}

\textbf{LLM:} Ich passe an:

\textbf{Kontext-Änderungen (Ψ):}
\begin{itemize}[nosep]
\item Geography: CH → \textbf{DE}
\item Regulatory: BVG → \textbf{BetrAVG}
\item Autonomy: 0.85 → \textbf{0.60} (DE niedriger)
\item Trust: 0.85 → \textbf{0.75}
\end{itemize}

\textbf{Parameter-Anpassungen:}
\begin{itemize}[nosep]
\item γ\_temporal: 0.20 → \textbf{0.15} (weniger Present Bias in DE)
\item β\_S: 0.20 → \textbf{0.25} (mehr Arbeitgeber-Einfluss)
\item Add: γ\_complexity = \textbf{-0.20} (BetrAVG komplexer als BVG)
\end{itemize}

Soll ich weitere Anpassungen vornehmen?

\textbf{User:} Nein, passt.

\hrule
\vspace{0.5em}

\textbf{LLM:} Hier ist Ihr angepasstes Modell:

\vspace{0.3em}
\fbox{\parbox{0.95\linewidth}{
\textbf{\large BetrAV Enrollment Modell (Deutschland, L0)}

\vspace{0.3em}
\textbf{Basis:} SEED-CH-01 (adaptiert)

\vspace{0.3em}
\textbf{━━━ MODELL-GLEICHUNG ━━━}
\begin{equation*}
P(\text{enroll}) = \frac{1}{1+\exp\bigl(-(\beta_0 + 0.35\,F + 0.25\,E + 0.25\,S + 0.30\,F \cdot E - 0.20\,\text{complexity})\bigr)}
\end{equation*}

\textbf{━━━ ANGEPASSTE VORHERSAGEN ━━━}

\textit{Bei Baseline von 50\% Enrollment (DE niedriger als CH):}
\begin{enumerate}[nosep]
\item \textbf{Framing ``Altersvorsorge'':} +6--10pp → 56--60\% \\
      \textit{(reduziert wg. niedrigerer Autonomie)}
\item \textbf{Arbeitgeber-Empfehlung:} +8--12pp → 58--62\% \\
      \textit{(verstärkt wg. höherem S in DE)}
\item \textbf{Vereinfachte Erklärung:} +5--8pp → 55--58\% \\
      \textit{(reduziert complexity barrier)}
\end{enumerate}

\textbf{━━━ DE-SPEZIFISCHE EMPFEHLUNG ━━━}

In Deutschland ist der \textbf{Arbeitgeber-Kanal} wichtiger als in der Schweiz.
Kombination aus Arbeitgeber-Empfehlung + Vereinfachung ergibt ca. +12--18pp.
}}

\vspace{0.3em}
\textit{Gesamtzeit: $<$5 Minuten | Basis: SEED-CH-01 v1.0}
\end{tcolorbox}


% =============================================================================
%                    PART 4: BACK MATTER
% =============================================================================

%%%%%%%%%%%%%%%%%%%%%%%%%%%%%%%%%%%%%%%%%%%%%%%%%%%%%%%%%%%%%%%%%%%%%%%%%%%%%%%
\section{Summary}
\label{sec:eee-summary}
%%%%%%%%%%%%%%%%%%%%%%%%%%%%%%%%%%%%%%%%%%%%%%%%%%%%%%%%%%%%%%%%%%%%%%%%%%%%%%%

\begin{tcolorbox}[colback=green!10!white, colframe=green!75!black,
    title=\textbf{Appendix DSN: Summary}]

\textbf{Core Question:} How do we systematically design testable models from EBF?

\textbf{Main Contribution:}
\begin{itemize}[nosep]
    \item Defines the Model Design Workflow with \textbf{Recursive 3+1 Architecture}
    \item Step 0: Workflow selection (Geführt / Schnell / Template / Custom)
    \item Steps 1--9: Full guided workflow for GEFÜHRT mode
    \item Introduces 3+1 Choice Architecture (Axiom MD-1)
    \item Links option selection to Model Registry (Axiom MD-2)
    \item \textbf{Provides LLM Integration for automated model design}
\end{itemize}

\textbf{The 4 Workflow Paths:}
\begin{itemize}[nosep]
    \item \textbf{GEFÜHRT:} Full 9-step workflow, 3+1 options at each step
    \item \textbf{SCHNELL:} 3 questions → \textbf{10C-complete} model via GGG Mappings 1--7
    \item \textbf{TEMPLATE:} Start from seed model (FFF), adapt differences
    \item \textbf{CUSTOM:} Full user control, no preset structure
\end{itemize}

\textbf{Key Axioms:}
\begin{itemize}[nosep]
    \item MD-1: 3+1 Choice Architecture (applies recursively to workflow selection)
    \item MD-2: Registry-Based Option Selection
    \item MD-3: Scope-Context Iteration
    \item MD-4: Option Selection Principles
\end{itemize}

\textbf{LLM Integration:}
\begin{itemize}[nosep]
    \item System prompt for Model Design Assistant (Section \ref{sec:eee-llm})
    \item Step 0 determines conversation path
    \item \textbf{Direct model output:} Gleichung, Vorhersagen, Interventionen
    \item GGG mapping table integration for auto-populated defaults
\end{itemize}

\textbf{Integration:} This appendix connects to FFF (Registry), GGG (Configurator),
all 8 CORE appendices (AAA, B, C, V, BBB, AU, AV, AW), and all DOMAIN/PREDICT appendices.

\end{tcolorbox}


%%%%%%%%%%%%%%%%%%%%%%%%%%%%%%%%%%%%%%%%%%%%%%%%%%%%%%%%%%%%%%%%%%%%%%%%%%%%%%%
\section{Glossary of Symbols}
\label{sec:eee-glossary}
%%%%%%%%%%%%%%%%%%%%%%%%%%%%%%%%%%%%%%%%%%%%%%%%%%%%%%%%%%%%%%%%%%%%%%%%%%%%%%%

\begin{tcolorbox}[colback=blue!5!white, colframe=blue!75!black]
\textbf{Central Reference:} For complete symbol definitions, see
\textbf{Appendix GLS} (\textit{Glossary and Notation}, \S\ref{app:glossary}).

This section lists only symbols introduced or specialized in this appendix.
\end{tcolorbox}

\vspace{0.5em}

\begin{table}[htbp]
\centering
\caption{Symbols Introduced in Appendix DSN}
\label{tab:eee-symbols}
\renewcommand{\arraystretch}{1.3}
\begin{tabular}{@{}clp{6cm}c@{}}
\toprule
\textbf{Symbol} & \textbf{Name} & \textbf{Definition} & \textbf{Also in G?} \\
\midrule
$E$ & Decision Point & A point in workflow requiring choice & NEW \\
$O_i$ & Option & Curated option $i \in \{1,2,3\}$ & NEW \\
$O_{\text{custom}}$ & Custom Option & Self-defined option & NEW \\
$\mathcal{R}$ & Registry & Set of validated models in FFF & NEW \\
$S$ & Scope & Behavior domain to be modeled & NEW \\
$M$ & Model & Complete behavioral model & \checkmark \\
$\Psi$ & Context & Context vector & \checkmark \\
$C$ & Utility Dimensions & FEPSDE dimensions & \checkmark \\
$\gamma$ & Complementarity & Interaction parameters & \checkmark \\
$\Theta$ & Parameters & Model parameters & \checkmark \\
\bottomrule
\end{tabular}
\end{table}


%%%%%%%%%%%%%%%%%%%%%%%%%%%%%%%%%%%%%%%%%%%%%%%%%%%%%%%%%%%%%%%%%%%%%%%%%%%%%%%

%%%%%%%%%%%%%%%%%%%%%%%%%%%%%%%%%%%%%%%%%%%%%%%%%%%%%%%%%%%%%%%%%%%%%%%%%%%%%%%
\section{Critical Foundations and Objections}
\label{sec:dsn-foundations}
%%%%%%%%%%%%%%%%%%%%%%%%%%%%%%%%%%%%%%%%%%%%%%%%%%%%%%%%%%%%%%%%%%%%%%%%%%%%%%%

\subsection{Objection 1: Scope Limitations}

\textbf{Objection:} The framework presented may have limited applicability
outside the specific domain contexts discussed.

\textbf{Response:} We acknowledge domain-specificity. The framework provides
general principles that should be calibrated to specific contexts. Cross-domain
validation remains an open research question.

\subsection{Objection 2: Measurement Challenges}

\textbf{Objection:} Key constructs may be difficult to measure reliably in
practice.

\textbf{Response:} We recommend using validated scales and instruments where
available, and developing domain-specific proxies where necessary. Sensitivity
analysis should assess robustness to measurement uncertainty.


\section{Open Issues and Research Agenda}
\label{sec:eee-open}
%%%%%%%%%%%%%%%%%%%%%%%%%%%%%%%%%%%%%%%%%%%%%%%%%%%%%%%%%%%%%%%%%%%%%%%%%%%%%%%

\begin{table}[htbp]
\centering
\caption{Research Roadmap for METHOD-DESIGN}
\label{tab:eee-roadmap}
\renewcommand{\arraystretch}{1.3}
\begin{tabular}{@{}clcl@{}}
\toprule
\textbf{\#} & \textbf{Issue} & \textbf{Priority} & \textbf{Status} \\
\midrule
1 & Similarity metric for Registry search & HIGH & Open \\
2 & Automation of 3+1 option generation & MEDIUM & \textbf{Implemented} (§\ref{sec:eee-llm}) \\
3 & Optimal iteration count for Scope-Context & LOW & Open \\
4 & Integration with AI-assisted design & MEDIUM & \textbf{Implemented} (§\ref{sec:eee-llm}) \\
5 & Cross-cultural validation of workflow & LOW & Planned \\
6 & LLM prompt optimization & MEDIUM & Open \\
\bottomrule
\end{tabular}
\end{table}


%%%%%%%%%%%%%%%%%%%%%%%%%%%%%%%%%%%%%%%%%%%%%%%%%%%%%%%%%%%%%%%%%%%%%%%%%%%%%%%
\section{References}
\label{sec:eee-references}
%%%%%%%%%%%%%%%%%%%%%%%%%%%%%%%%%%%%%%%%%%%%%%%%%%%%%%%%%%%%%%%%%%%%%%%%%%%%%%%

\begin{tcolorbox}[colback=blue!5!white, colframe=blue!75!black]
\textbf{Master Bibliography:} For complete reference details, see
\texttt{bcm2\_master\_references.bib} in the repository.

This section lists appendix-specific references and data sources.
\end{tcolorbox}

\vspace{0.5em}

\subsection{Key References for This Appendix}

\begin{itemize}[nosep]
    \item \textbf{Choice Architecture:} Thaler \& Sunstein (2008), Johnson et al. (2012)
    \item \textbf{Decision Fatigue:} Baumeister et al. (2008), Vohs et al. (2014)
    \item \textbf{Model Design:} Box (1976), Friedman (1953)
\end{itemize}

\nocite{bcm_master}


%%%%%%%%%%%%%%%%%%%%%%%%%%%%%%%%%%%%%%%%%%%%%%%%%%%%%%%%%%%%%%%%%%%%%%%%%%%%%%%
% CROSS-REFERENCE FOOTER (Required)
%%%%%%%%%%%%%%%%%%%%%%%%%%%%%%%%%%%%%%%%%%%%%%%%%%%%%%%%%%%%%%%%%%%%%%%%%%%%%%%

\vspace{1em}
\hrule
\vspace{0.5em}

\noindent\textit{Cross-references:}
Appendix GLS (Central Glossary),
Appendix REG (METHOD-REGISTRY),
Appendix CFG (METHOD-CONFIG),
Appendix B (CORE-HOW),
Appendix C (CORE-WHAT),
Appendix V (CORE-WHEN),
Appendix BBB (CORE-WHERE),
Appendix CAL (METHOD-LLMMC).

\noindent\textit{Master Bibliography:} \texttt{bcm2\_master\_references.bib}


% =============================================================================
% END OF APPENDIX EEE
% =============================================================================
